\begin{figure}[htbp]
	\centering
	\begin{tikzpicture}[
		>=Stealth, line width=0.8pt, scale=1.0,
		asta/.style={line width=1.6pt},
		fz/.style={-{Stealth[length=2.2mm]}, line width=0.9pt},
		et/.style={font=\footnotesize},
		nd/.style={circle, fill=white, draw=black,
			inner sep=0pt, minimum size=5pt}]
		
		% ═══════════ PANNELLO (a): nodo non isolato ═══════════
		\begin{scope}[xshift=0cm]
			\def\R{1.25}
			\foreach \a/\nome in {90/1, -30/2, -150/3}{
				\draw[asta] (0,0)--(\a:\R);
				\node[et] at (\a:{\R+0.30}) {$e_{\nome}$};}
			% cerniera tripla: un solo pallino
			\node[nd, minimum size=7pt] at (0,0) {};
			\node[et, above left=1pt] at (0,0) {$i$};
			\node[et, align=center] at (0,-2.20)
			{cerniera tripla\\
				$2(n_i-1) = 4$ incognite\\
				nessuna equazione propria};
		\end{scope}
		
		% ═══════════ PANNELLO (b): nodo isolato ═══════════
		\begin{scope}[xshift=6.4cm]
			\def\R{1.25}
			\def\g{0.42}   % ampiezza del distacco
			\foreach \a/\nome in {90/1, -30/2, -150/3}{
				\draw[asta] (\a:\g)--(\a:\R);
				\node[et] at (\a:{\R+0.30}) {$e_{\nome}$};
				% forza trasmessa al nodo: due componenti
				\draw[fz] (\a:\g)--($(\a:\g)+(\a:-0.32)$);}
			\node[nd, minimum size=7pt] at (0,0) {};
			\node[et, above left=1pt] at (0,0) {$i$};
			\node[et, align=center] at (0,-2.20)
			{nodo isolato\\
				$2n_i = 6$ incognite\\
				$+\,2$ equazioni di equilibrio};
		\end{scope}
		
		% ═══════════ bilancio ═══════════
		\node[et] at (3.2,-3.10)
		{$+2$ incognite \quad contro \quad $+2$ equazioni};
	\end{tikzpicture}
	\caption{Le due letture di un nodo a cerniera in cui
		convergono tre travi. A sinistra il nodo è un vincolo
		interno e le sue incognite statiche sono quattro,
		secondo la \EQUref{eq:molt_nodo}; non porta equazioni
		proprie. A destra il nodo è isolato come corpo libero:
		ciascuna trave scambia con esso una forza di due
		componenti, per un totale di sei incognite, e il nodo
		porta le due equazioni
		\eqref{eq:equilibrio_nodo_cerniera}. Le due incognite
		in più sono compensate dalle due equazioni in più.}
	\label{fig:partizione_nodo}
\end{figure}
